%% 附录

\begin{thesisappendix}
\section{数据集及构建补充材料}
\label{app:dataset-construction}
\label{app:toolchain-entry}

本附录补充 UAV-DualCog 测试基准的数据规模、任务结构和构建流程。正文第3章和第4章分别说明了双认知任务设计与四阶段工具链，本附录进一步给出补充材料中用于复核的结构化信息，使数据集构建过程、任务协议和实现入口之间形成可追溯关系。

\begin{table}[htbp]
    \centering
    \caption{UAV-DualCog 测试基准发布规模补充统计}
    \label{tab:app-dataset-scale}
    \small
    \begin{tabularx}{\textwidth}{@{}X>{\centering\arraybackslash}p{3.0cm}X@{}}
        \toprule
        统计项 & 数值 & 说明 \\
        \midrule
        测试场景数 & 12 & 来自与训练场景隔离的 AerialVLN 仿真环境 \\
        有效地标数 & 512 & 经过多视角复核和语义标注后的发布地标 \\
        图像任务样本数 & 4096 & 覆盖四类图像任务与两条双认知能力主线 \\
        视频任务样本数 & 2048 & 覆盖飞行行为识别与地标可见性推理 \\
        图像发布分辨率 & DCI 4K / JPEG & 保留高分辨率地标细节，评测时可按模型能力压缩输入 \\
        视频发布规格 & 1080P / 10 FPS / MP4 & 支持行为识别、可见性计数和时间区间定位 \\
        \bottomrule
    \end{tabularx}
\end{table}

\begin{table}[htbp]
    \centering
    \caption{UAV-DualCog 六类任务与证据形式}
    \label{tab:app-task-matrix}
    \small
    \setlength{\tabcolsep}{4pt}
    \begin{tabularx}{\textwidth}{@{}p{3.0cm}p{1.8cm}p{2.2cm}X p{2.8cm}@{}}
        \toprule
        任务类型 & 媒介 & 认知分支 & 输入形式 & 输出与证据 \\
        \midrule
        地标参照位置推理 & 图像 & 自我状态认知 & 地标参考图 + 当前观测图 & 选项 + 目标框 \\
        未来观察预测 & 图像 & 自我状态认知 & 地标参考图 + 候选未来视图 & 选项 + 目标框 \\
        自我参照位置推理 & 图像 & 环境情况认知 & 当前观测图 + 地标描述 & 选项 + 目标框 \\
        地标驱动动作决策 & 图像 & 环境情况认知 & 当前观测图 + 地标描述 & 选项 + 目标框 \\
        飞行行为识别与时间定位 & 视频 & 自我状态认知 & 第一视角飞行视频 & 行为选项 + 时间区间 \\
        地标可见次数与时间定位 & 视频 & 环境情况认知 & 参考图 + 第一视角视频 & 可见次数 + 时间区间 \\
        \bottomrule
    \end{tabularx}
\end{table}

\begin{table}[htbp]
    \centering
    \caption{测试基准构建工具链主要入口}
    \label{tab:app-toolchain}
    \scriptsize
    \begin{tabularx}{\textwidth}{@{}p{1.7cm}p{3.8cm}X@{}}
        \toprule
        阶段 & 主要入口 & 作用说明 \\
        \midrule
        阶段1 & \makecell[l]{\texttt{stage1\_collect\_}\\\texttt{pcd.py}} & 调用 AerialVLN 模拟器进行场景扫描，采集语义分割、位姿和点云信息，并融合为场景级几何基础。 \\
        阶段2 & \makecell[l]{\texttt{stage2\_landmark\_}\\\texttt{label.py}} & 基于实例点云生成地标候选，支持人工复核、语义标注和有效地标清单导出。 \\
        阶段3 & \makecell[l]{\texttt{stage3\_generate\_}\\\texttt{traj.py}\\\texttt{stage3\_task\_}\\\texttt{suite.py}} & 生成飞行轨迹、视频资产、复合行为和原子行为时间监督，并导出视频任务。 \\
        阶段4 & \makecell[l]{\texttt{stage4\_qa\_generate\_}\\\texttt{and\_eval.py}} & 围绕地标中心参考视图和无人机观测图生成图像问答任务，并导出目标框监督。 \\
        全流程 & \makecell[l]{\texttt{task\_pipeline.py}\\\texttt{pipeline\_common.py}} & 组织阶段间输入输出、路径管理、任务配置和批处理运行。 \\
        \bottomrule
    \end{tabularx}
\end{table}

\subsection{四阶段构建流程伪代码}

\begin{verbatim}
Input: AerialVLN scenes and task configuration
Output: UAV-DualCog benchmark with image/video evidence labels

for each selected scene do
    Stage 1: scan scene and collect RGB / semantics / LiDAR / poses
    fuse observations into a scene-level semantic point cloud

    Stage 2: cluster landmark candidates from semantic point cloud
    render multi-view landmark candidates
    conduct human review and semantic labeling
    export validated landmark assets

    Stage 3: bind landmarks with flight behavior templates
    generate trajectories and first-person videos
    record temporal evidence for actions and landmark visibility
    export self-aware and environment-aware video tasks

    Stage 4: sample landmark-centered reference view
    sample UAV observation view or future-view candidates
    generate image QA tasks and normalized bbox evidence
end for
\end{verbatim}

阶段3的关键在于为同一条轨迹同时建立高层行为语义与底层时间证据。阶段4的关键在于同步生成选项答案和目标框证据，而不是先生成问题再事后补充定位框。因此，UAV-DualCog 的任务构建天然具备“语义答案 + 显式证据”的联合评测属性。
\end{thesisappendix}

\begin{thesisappendix}
\section{评测协议补充材料}
\label{app:evaluation-protocol}
\label{app:prompt-protocol}

UAV-DualCog 的评测协议由系统提示词、用户提示词模板、结构化输出协议和指标解析规则共同组成。系统提示词统一限定模型扮演无人机智能体、明确输入媒介和禁止额外解释；用户提示词注入地标描述、候选选项、参考图像、查询图像或视频语境；输出协议要求模型返回严格 JSON，便于解析器稳定计算语义指标和证据指标。

\begin{table}[htbp]
    \centering
    \caption{评测输出字段与指标映射}
    \label{tab:app-output-metric-map}
    \footnotesize
    \begin{tabularx}{\textwidth}{@{}p{3.0cm}p{3.4cm}X@{}}
        \toprule
        输出字段 & 适用任务 & 驱动指标 \\
        \midrule
        \makecell[l]{\texttt{answer\_}\\\texttt{option\_id}} & 图像任务 & Option Accuracy \\
        \texttt{bbox\_xyxy\_norm} & 图像任务 & BBox Acc@0.5、BBox Mean IoU \\
        \makecell[l]{\texttt{answers[]}\\\texttt{.intervals\_sec}} & 自我状态认知视频任务 & 行为识别主指标、Temporal F1@0.5、Mean tIoU \\
        \texttt{visible\_count} & 环境情况认知视频任务 & Count Accuracy、Count Within1 \\
        \makecell[l]{\texttt{visible\_}\\\texttt{intervals\_sec}} & 环境情况认知视频任务 & Segment F1@0.5、Mean tIoU \\
        \bottomrule
    \end{tabularx}
\end{table}

\begin{table}[htbp]
    \centering
    \caption{六类任务提示词协议摘要}
    \label{tab:app-prompt-protocol}
    \footnotesize
    \setlength{\tabcolsep}{4pt}
    \begin{tabularx}{\textwidth}{@{}p{2.8cm}p{2.5cm}X p{2.8cm}@{}}
        \toprule
        任务 & 输入媒介 & 核心提示信息 & 必要输出字段 \\
        \midrule
        地标参照位置推理 & 参考图像 + 查询图像 & 地标中心参考视图、当前无人机观测图、相对位置选项 & \makecell[l]{\texttt{answer\_option\_id}\\\texttt{bbox\_xyxy\_norm}} \\
        未来观察预测 & 参考图像 + 候选图像 & 当前观察、轨道运动描述、多张候选未来观察图 & \makecell[l]{\texttt{answer\_option\_id}\\\texttt{bbox\_xyxy\_norm}} \\
        自我参照位置推理 & 查询图像 & 当前第一视角观察、地标描述、方位选项 & \makecell[l]{\texttt{answer\_option\_id}\\\texttt{bbox\_xyxy\_norm}} \\
        地标驱动动作决策 & 查询图像 & 当前第一视角观察、目标地标描述、动作选项 & \makecell[l]{\texttt{answer\_option\_id}\\\texttt{bbox\_xyxy\_norm}} \\
        飞行行为识别与时间定位 & 视频 & 第一视角飞行视频、复合或原子行为候选选项 & \makecell[l]{\texttt{answers}\\\texttt{intervals\_sec}} \\
        地标可见次数与时间定位 & 视频 + 参考图像 & 目标地标参考图、第一视角飞行视频、可见性问题 & \makecell[l]{\texttt{visible\_count}\\\texttt{visible\_}\\\texttt{intervals\_sec}} \\
        \bottomrule
    \end{tabularx}
\end{table}

\subsection{图像任务完整中文提示词模板}

\paragraph{地标参照位置推理。}
系统提示词要求模型作为 UAV agent，对比包含目标地标的参考视图和当前第一视角视图，根据参考图中可见立面定义的地标中心坐标系，判断当前无人机相对地标的位置，并同时定位查询图像中的目标地标。输出必须是严格 JSON：
\begin{verbatim}
{"answer_option_id":"A","bbox_xyxy_norm":[0.1,0.2,0.3,0.4]}
\end{verbatim}
用户提示词模板为：图像1显示“参考立面变量”对应立面的“地标描述变量”，该立面位于地标中心坐标系中。基于这一参考，图像2中无人机相对于该地标处于什么位置？请选择一个选项，并返回图像2中该地标的归一化目标框。

\paragraph{未来观察预测。}
系统提示词要求模型根据当前包含地标的参考视图和给定轨道运动，选择运动后最可能得到的未来观察图，并在所选候选图像中定位目标地标。用户提示词模板为：图像1显示“参考立面变量”对应立面的“地标描述变量”。执行“运动实例变量”后，哪一张候选图像是正确的未来视图？请选择一个选项，并返回所选图像中该地标的归一化目标框。

\paragraph{自我参照位置推理。}
系统提示词要求模型从无人机当前第一视角出发，判断目标地标相对无人机前进方向的位置，并输出该地标的空间证据。用户提示词模板为：图像1显示“参考立面变量”对应立面的“地标描述变量”。在图像2中，该地标相对于无人机前进方向位于哪里？请选择一个选项，并返回图像2中该地标的归一化目标框。

\paragraph{地标驱动动作决策。}
系统提示词要求模型判断无人机为了接近目标地标应采取的飞行方向，并用目标框说明决策所依据的地标位置。用户提示词模板为：图像1显示“参考立面变量”对应立面的“地标描述变量”。在图像2中，为了接近该地标，无人机应该向哪个方向飞行？请选择一个选项，并返回图像2中该地标的归一化目标框。

\subsection{视频任务完整中文提示词模板}

\paragraph{复合飞行行为识别。}
系统提示词要求模型观看无人机第一视角飞行视频，识别视频中曾经发生过的飞行动作，可以选择多个选项，并返回每个选项对应的时间区间。输出必须是严格 JSON：
\begin{verbatim}
{"answers":[{"option_id":"A","intervals_sec":[[0.0,1.2]]}]}
\end{verbatim}
用户提示词模板为：在该第一视角飞行视频中，无人机曾经执行过以下哪些动作？请选择一个或多个选项，并返回对应的秒级时间区间。

\paragraph{原子飞行行为识别。}
系统提示词要求模型识别无人机在第一视角视频中执行的原子动作，并返回匹配时间区间。用户提示词模板为：在该第一视角飞行视频中，无人机执行了以下哪一个动作？请选择一个选项，并返回对应的秒级时间区间；若多个候选动作确实出现，则需要全部返回。

\paragraph{地标可见次数与时间定位。}
系统提示词要求模型结合目标地标参考图和第一视角飞行视频，判断该地标在视频中分几次可见，并恢复完整可见时间区间。输出必须是严格 JSON：
\begin{verbatim}
{"visible_count":2,"visible_intervals_sec":[[0.0,9.1],[17.7,26.8]]}
\end{verbatim}
用户提示词模板为：参考图显示“地标描述变量”。该地标在无人机第一视角飞行视频中出现了几次？请返回完整可见时间区间，单位为秒。

\subsection{解析与有效性规则}

所有任务均要求严格 JSON 输出。包含 Markdown 代码块、额外自然语言解释、字段缺失或字段类型错误的结果会被视为格式无效。图像任务只有在选项与目标框均可解析时才进入完整评分；视频任务则将语义识别或可见次数与时间区间质量分开计算，从而避免把“识别正确但证据错误”的输出误判为完整成功。
\end{thesisappendix}

\begin{thesisappendix}
\section{训练集及构建补充材料}
\label{app:train-dataset-construction}

UAV-DualCog-Train-V1 面向第6章能力优化实验构建，仅使用与测试集场景隔离的训练场景。其目标不是扩大测试集，而是将第五章暴露出的“图像语义判断与空间证据脱节”问题转化为可监督、可奖励的训练记录。

\begin{table}[htbp]
    \centering
    \caption{UAV-DualCog-Train-V1 训练数据字段设计}
    \label{tab:app-train-fields}
    \small
    \begin{tabularx}{\textwidth}{@{}p{3.1cm}p{3.0cm}X@{}}
        \toprule
        字段 & 类型 & 作用 \\
        \midrule
        \texttt{prompt} & 对话文本 & 保存与评测协议一致的系统提示词和用户提示词 \\
        \texttt{completion} & JSON 字符串 & 保存监督目标，包括 \texttt{answer\_option\_id} 和 \texttt{bbox\_xyxy\_norm} \\
        \texttt{images} & 图像路径列表 & 按评测输入顺序保存真实图像路径 \\
        \texttt{task\_type} & 字符串 & 区分四类图像训练任务 \\
        \texttt{difficulty} & 字符串 & 区分两档问题难度 \\
        \makecell[l]{\texttt{answer\_bbox}\\\texttt{\_xyxy\_1000}} & 数值列表 & 供奖励函数复用的目标框辅助字段 \\
        \bottomrule
    \end{tabularx}
\end{table}

\begin{table}[htbp]
    \centering
    \caption{UAV-DualCog-Train-V1 数据分布与完整性补充统计}
    \label{tab:app-train-stat}
    \footnotesize
    \setlength{\tabcolsep}{4pt}
    \begin{tabularx}{\textwidth}{@{}>{\raggedright\arraybackslash}X>{\centering\arraybackslash}p{2.2cm}>{\raggedright\arraybackslash}X>{\centering\arraybackslash}p{2.2cm}@{}}
        \toprule
        统计项 & 数值 & 统计项 & 数值 \\
        \midrule
        有效训练地标数 & 234 & 全量样本数 & 11232 \\
        训练集样本数 & 10104 & 验证集样本数 & 1128 \\
        四类任务样本数 & 每类 2808 & 双认知分组样本数 & 各 5616 \\
        两档难度样本数 & 各 5616 & 每个任务-难度组合 & 1404 \\
        非法目标框数量 & 0 & 图像路径存在率 & 100\% \\
        目标框面积均值 & 0.0645 & 训练/验证比例 & 89.96\% : 10.04\% \\
        \bottomrule
    \end{tabularx}
\end{table}

\subsection{训练数据准备伪代码}

\begin{verbatim}
Input: stage4 manifests from six training scenes
Output: all / train / valid JSONL and integrity report

for each scene manifest do
    load image-task samples
    build benchmark-aligned system/user prompt
    collect real image paths in benchmark input order
    build assistant completion as JSON:
        {"answer_option_id": "...",
         "bbox_xyxy_norm": [...]}
    cache normalized bbox and bbox_1000 for reward computation
end for

bucket samples by (task_type, difficulty)
sort each bucket by sample id
assign every 10th sample to valid split
assign remaining samples to train split
dump JSONL files and integrity report
\end{verbatim}

该构建方式保证训练样本在任务类型、难度等级和双认知分组上保持均衡，同时保证训练集与测试集在场景层面隔离。训练集没有引入视频样本，因此第6章的视频结果被解释为跨媒介迁移检验，而不是视频任务的直接监督收益。
\end{thesisappendix}

\begin{thesisappendix}
\section{能力优化实现补充材料}
\label{app:optimization-entry}

第6章正文保留能力优化方法、关键参数、训练完整性和最终任务级评测结果。本附录补充实现入口、训练流程伪代码、完整性核查逻辑和奖励函数关键片段，用于说明方法已经落实为可执行训练链路。

\begin{table}[htbp]
    \centering
    \caption{能力优化主要入口与作用说明}
    \label{tab:app-optimization-entry}
    \scriptsize
    \begin{tabularx}{\textwidth}{@{}p{2.6cm}p{4.1cm}X@{}}
        \toprule
        脚本入口 & 主要输入输出 & 作用说明 \\
        \midrule
        \makecell[l]{\texttt{prepare\_train}\\\texttt{\_dataset.py}} & 输入 Stage 4 manifest；输出 train/valid/all JSONL 与完整性报告 & 将 UAV-DualCog-Train-V1 图像任务转换为统一多模态训练记录。 \\
        \makecell[l]{\texttt{run\_sft}\\\texttt{\_lora.py}} & 输入 train/valid JSONL 与 Qwen 3.5 4B；输出 LoRA adapter、训练指标与合并模型 & 实现多模态监督微调，使真实图像参与训练并学习结构化 JSON 输出。 \\
        \makecell[l]{\texttt{run\_grpo}\\\texttt{\_light.py}} & 输入训练 JSONL 与 SFT 检查点；输出 GRPO adapter、奖励指标与合并模型 & 实现联合 GRPO，在同一奖励函数中优化格式、答案和目标框质量。 \\
        \makecell[l]{\texttt{run\_full}\\\texttt{\_pipeline.sh}} & 管理数据快照、日志、检查点和两阶段调度 & 串联数据准备、SFT 与 Joint GRPO，形成单机多卡可复现训练流程。 \\
        \makecell[l]{\texttt{check\_training}\\\texttt{\_run.py}} & 输入运行目录；输出完整性报告 & 核查计划文件、日志、checkpoint、adapter 和合并模型产物是否完整。 \\
        \makecell[l]{\texttt{evaluate}\\\texttt{\_transfer.py}} & 输入 Base/GRPO 权重与 benchmark 场景；输出评测结果与增益表 & 触发 Stage 3/4 评测并生成最终优化模型对比结果。 \\
        \bottomrule
    \end{tabularx}
\end{table}

\subsection{两阶段训练流程伪代码}

\begin{verbatim}
Input: UAV-DualCog-Train-V1 JSONL, Qwen 3.5 4B base model
Output: SFT checkpoint and final GRPO checkpoint

Stage 1: SFT
    load train / valid JSONL
    load real RGB images on demand
    encode text and images with AutoProcessor
    train LoRA adapters with completion-only loss
    export adapter and merged model

Stage 2: Joint GRPO
    initialize from SFT checkpoint
    load real RGB images on demand
    generate JSON completions under multimodal model
    compute format reward + answer reward + bbox reward
    update LoRA adapters with group-relative optimization
    export final adapter and merged model
\end{verbatim}

\subsection{正式训练完整性核查逻辑}

正式训练前进行了小样本链路验证：SFT 验证使用 8 条样本和 1 个优化步检查真实图像加载、前向反向传播和模型导出；GRPO 验证使用 4 条样本和 1 个优化步检查奖励计算、组相对优化和合并模型导出。该验证仅用于确认工程链路可执行，不作为能力提升结论。

\begin{verbatim}
Input: training run directory
Output: integrity JSON / Markdown report

check run root exists
for each stage in [01_sft_full, 02_grpo_joint] do
    verify plan file, metrics file, log file exist
    verify adapter and merged_model exist
    count checkpoints and locate latest step
    parse trainer_state tail summary
end for

if both stages complete and merged models exist then
    mark overall_complete = true
\end{verbatim}

\subsection{奖励函数关键代码片段}
\label{app:reward-code}

下列代码片段节选自奖励实现，展示本文如何优先解析 JSON 结构化输出，并将格式合法性、答案正确性和目标框一致性组合为统一奖励。
\begin{verbatim}
def reward_format(text: str) -> float:
    return 1.0 if (
        parse_answer(text) is not None
        and parse_bbox_norm(text) is not None
    ) else 0.0

def reward_mcq(pred_text: str, gt_answer: str) -> float:
    pred = parse_answer(pred_text)
    if pred is None:
        return 0.0
    return 1.0 if pred == str(gt_answer).strip().upper() else 0.0

def reward_bbox(pred_text: str, gt_bbox_norm):
    pred_bbox = parse_bbox_norm(pred_text)
    if pred_bbox is None:
        return 0.0
    giou = giou_norm(pred_bbox, gt_bbox_norm)
    return max(0.0, min(1.0, (giou + 1.0) * 0.5))

def total_reward(pred_text, gt_answer, gt_bbox_norm, weights):
    return (
        weights.w_format * reward_format(pred_text)
        + weights.w_mcq * reward_mcq(pred_text, gt_answer)
        + weights.w_bbox * reward_bbox(pred_text, gt_bbox_norm)
    )
\end{verbatim}

该实现与第6章中的奖励公式一一对应，保证能力优化目标并非停留在方法设想层面，而是直接进入了最终 GRPO 训练与任务级评测流程。
\end{thesisappendix}
